% Lösungen Ableitung und Änderungsrate (zusammenbau v0.4, Heft)
\einheitenkopf{Ableitung und Änderungsrate}

\erg{1}{a) Zeitraum – Differenzenquotient, Sekante: „in den ersten fünf Stunden“ nennt ein Intervall \quad b) $\frac{300 - 120}{3} = 60$ Besucher je Stunde \quad c) $A$: $\frac{A(10) - A(0)}{10} \approx -4{,}42$; $B$: $\frac{60 - 90}{10} = -3$; der Betrag ist bei $A$ größer}
\erg{2}{a) $a(1) = 100$, $a(3) = 900$; $\frac{900 - 100}{2} = 400$ Bewertungen je Stunde \quad b) $h(5) \approx 63{,}36$, $h(0) = 40$; $\frac{h(5) - h(0)}{5} \approx 4{,}67$ Prozentpunkte je Tag \quad c) $k(16) = 32{,}6$, $k(10) = 5$; $\frac{k(16) - k(10)}{6} = 4{,}6$ mmol/l je km/h \quad d) I: $2 - 1{,}1 = 0{,}9$ Liter – so viel Luft wird beim Einatmen zwischen $1{,}5$ und $3$ Sekunden aufgenommen; II: $0{,}9 : 1{,}5 = 0{,}6$ Liter je Sekunde – die mittlere Einatemrate in diesem Zeitraum \quad e) $f(1) \approx 2{,}011$, $f(5) \approx 0{,}947$; $m = \frac{f(5) - f(1)}{4} \approx -0{,}266$; $y = -0{,}266 \cdot (x - 1) + 2{,}011 \approx -0{,}266x + 2{,}277$ \quad f) $f$: $\frac{f(5) - f(0)}{5} \approx -0{,}211$; $h$: $\frac{0{,}6 - 1{,}8}{5} = -0{,}24$; der Betrag der mittleren Steigung ist bei $h$ größer ($0{,}24 > 0{,}211$)}
\erg{3}{a) Länge $4 \cdot 3{,}2 + 2\pi \cdot 0{,}5 \approx 15{,}94$ m; Zeit $\frac{15{,}94}{4} \approx 3{,}99$ min \quad b) $\frac{k(c) - k(0)}{c} = 0{,}1c^2 - 1{,}5c = -5$, also $c^2 - 15c + 50 = 0$ und $c = 5$ oder $c = 10$; kleinere Lösung $c = 5$: Die linke Seite ist die mittlere Abbaurate über $[0; c]$; in den ersten fünf Jahren nimmt der Bestand im Mittel um $5$ Mio. Tonnen je Jahr ab. \quad c) Falsch: $K(2) - K(1) = 10$, aber $K(3) - K(2) = 4$ – die Kosten des zusätzlichen Kubikmeters nehmen zunächst ab, obwohl $K$ steigt. \quad d) Falsch: $K(2) - K(1) = 15{,}5$, aber $K(3) - K(2) = 9{,}5$ – die Kosten des zusätzlichen Kubikmeters nehmen zunächst ab, obwohl $K$ steigt. \quad e) $g(x + 1) - g(x) = -500\mathrm{e}^{-0{,}4x} \cdot \mathrm{e}^{-0{,}4} + 500\mathrm{e}^{-0{,}4x} = 500(1 - \mathrm{e}^{-0{,}4})\mathrm{e}^{-0{,}4x}$ (Intervalllänge $1$); $500(1 - \mathrm{e}^{-0{,}4}) \approx 164{,}8$; $\mathrm{e}^{-0{,}4x} = \frac{80}{164{,}8}$ liefert $x \approx 1{,}81$ h, also nach etwa 1 h 48 min \quad f) $g(x + 1) - g(x) = -400\mathrm{e}^{-0{,}3x} \cdot \mathrm{e}^{-0{,}3} + 400\mathrm{e}^{-0{,}3x} = 400(1 - \mathrm{e}^{-0{,}3})\mathrm{e}^{-0{,}3x}$ (Intervalllänge $1$); $400(1 - \mathrm{e}^{-0{,}3}) \approx 103{,}7$; $\mathrm{e}^{-0{,}3x} = \frac{60}{103{,}7}$ liefert $x \approx 1{,}82$ h, also nach etwa 1 h 49 min}
\erg{4}{a) Ableitungswert – Steigung oder Rate: nach 4 Jahren wächst der Baum momentan mit $35$ cm je Jahr \quad b) $h'(t) = -10t + 20$, $h'(1) = 10$ m je s – die momentane Änderungsrate bei $t = 1$ \quad c) Um 14:30 Uhr steigt der Wasserstand, um 20:00 Uhr sinkt er.}
\erg{5}{a) Steigungsdreieck: 2 nach rechts, 6 nach oben, also Steigung $3$; die Ableitung einer linearen Funktion ist an jeder Stelle ihre Steigung, also $g'(3) = 3$ \quad b) $g'(x) = 3\mathrm{e}^x$, $g'(0) = 3$; Tangente im Punkt $(0 | 0)$ mit Steigungsdreieck 1 nach rechts, 3 nach oben \quad c) $g'(x) = 0{,}5\mathrm{e}^x$, $g'(0) = 0{,}5$; Tangente durch $(0 | 1{,}5)$, Steigungsdreieck 2 nach rechts, 1 nach oben \quad d) $(3 | 180)$: Drei Stunden nach Beobachtungsbeginn nimmt der Tankinhalt momentan mit $180$ kg je Stunde zu – $180$ ist eine Rate, kein Tankinhalt. \quad e) $f'(x) = 0{,}3x^2 - 2x + 2$, $f'(0) = 2$; $p'(x) = -2x + 3$, $p'(0{,}5) = 2$ – gleiche Steigung: die Tangenten an beide Begrenzungslinien in diesen Punkten sind parallel, die Schaufel läuft dort ohne Knick weiter.}
\erg{6}{a) Falsch: $u'(x) = 0{,}6x^2$, $v'(x) = 2x - 0{,}3x^2$; für $x = 2{,}5$ ist $u'(2{,}5) = 3{,}75$ und $v'(2{,}5) = 3{,}125$, also ist dort die Steigung von $u$ größer – eine Stelle genügt als Gegenbeispiel. \quad b) Falsch: $u'(x) = x$, $v'(x) = 3 - 0{,}5x$; für $x = 3$ ist $u'(3) = 3$ und $v'(3) = 1{,}5$ – dort ist $u$ steiler, die Allaussage ist widerlegt (gleich steil bei $x = 2$).}
\erg{7}{a) Sekantensteigung – mittlere Änderungsrate: Differenzenquotient über das Intervall von $2$ bis $5$ \quad b) $m = 5$; $4{,}5$; $4{,}1$; $4{,}01$ – die Steigungen nähern sich $4$: $f'(2) = 4$ \quad c) Mittlere Steigung $m = \frac{g(2) - g(-2)}{4} = \frac{2 - 4}{4} = -0{,}5$; Tangente in $W(0 | 3)$: Steigungsdreieck 2 nach rechts, 3 nach oben, Steigung $m_W = 1{,}5$ (positiv – mittlere Steigung und Wendetangente haben verschiedene Vorzeichen)}

7c)

\begin{ksys}[xmin=-3,xmax=3,ymin=-1,ymax=7,klein]\funktionab{-0.5*\x^3+1.5*\x+3}{g}{-2.4}{2.4}\tangentean{-0.5*\x^3+1.5*\x+3}{0}{t}\end{ksys}

\erg{8}{a) mittlere Rate $\frac{T(40) - T(0)}{40} \approx -0{,}249$ °C je min, momentane Rate $T'(40) \approx -0{,}225$ °C je min; Abweichung $\frac{0{,}249 - 0{,}225}{0{,}225} \approx 0{,}107 = 10{,}7\,\%$ – ja, mehr als $10\,\%$ (auf die mittlere Rate bezogen wären es nur $9{,}7\,\%$ – falsche Bezugsgröße, falsches Urteil) \quad b) mittlere Rate $\frac{B(2) - B(0)}{2} \approx 82{,}2$ Tausend je h, momentane Rate $B'(2) \approx 109{,}3$ Tausend je h; Abweichung $\frac{109{,}3 - 82{,}2}{109{,}3} \approx 0{,}248 = 24{,}8\,\%$ – nein, nicht mehr als $25\,\%$ (auf die mittlere Rate bezogen wären es $33\,\%$ – falsche Bezugsgröße, falsches Urteil) \quad c) Sekantensteigung $\frac{18 - 0}{3} = 6$; $f'(x_0) = 3x_0^2 - 3 = 6$ ⇔ $x_0^2 = 3$ ⇔ $x_0 = \sqrt{3} \approx 1{,}73$ (auch $-\sqrt{3}$); an dieser Stelle ist die Tangente parallel zur Sekante durch die Graphenpunkte bei $0$ und $3$ \quad d) Die Sekante durch die Schnittpunkte ist $g$ selbst, die mittlere Änderungsrate also $m = 0{,}5$ (keine Rechnung nötig); $f'(x) = \frac{1}{2\sqrt{x}} = 0{,}5$ ⇔ $\sqrt{x} = 1$ ⇔ $x = 1$}
\erg{9}{a) Bestand (Höhe, Anzahl, Menge) – größte Rate bei den Nullstellen der zweiten Ableitung: die Zuflussrate ist $f'$, ihr Maximum liegt bei $f'' = 0$ \quad b) $r'(t) = 0$ ⇔ $t = 6$ (der e-Faktor hat keine Nullstelle); nach 6 Stunden ist die Zuflussrate am größten \quad c) $T'(t) = -0{,}3t^2 + 3{,}6t$, $T''(t) = -0{,}6t + 3{,}6 = 0$ ⇔ $t = 6$; $T'(6) = 10{,}8$ °C je h}
\erg{10}{a) $f'(x) = 1{,}5x^2 - 3$, $f''(x) = 3x = 0$ ⇔ $x = 0$; kleinste Steigung $f'(0) = -3$ (nicht das Minimum von $f$ suchen) \quad b) $p''(x) = (0{,}32x - 1{,}44) \cdot \mathrm{e}^{-0{,}4x} = 0$ ⇔ $x = 4{,}5$; dort ist $p'(4{,}5) \approx -0{,}331 < -0{,}33$ (die mittlere Steigung im Intervall, etwa $-0{,}26$, genügt nicht als Beleg) \quad c) $w''(t) = -6t + 24 = 0$ ⇔ $t = 4$; $w'(4) = 48$ cm je Woche (nicht $t = 4$ als Ergebnis angeben) \quad d) $k''(x) = 0{,}3 \cdot (1 - 0{,}1x) \cdot \mathrm{e}^{-0{,}1x} = 0$ ⇔ $x = 10$; $k'(10) = 3 \cdot \mathrm{e}^{-1} \approx 1{,}10$ Mio. t je Jahr (nicht $k'(x) = 0$ lösen) \quad e) $h'(t) = 23{,}4$ ⇔ $t^2 - 20t + 36 = 0$ ⇔ $t = 2$ oder $t = 18$; $t = 18$ liegt außerhalb der Phase; $h'(0) = 45$: die Geschwindigkeit ist von $t = 0$ bis $t = 2$ mindestens $23{,}4$ m je min, der Zeitraum ist $2$ Minuten lang \quad f) $f'(t) = 9$ ⇔ $t^2 - 8t + 12 = 0$ ⇔ $t = 2$ oder $t = 6$; beide in der Phase: von $t = 2$ bis $t = 6$, die Länge des Zeitraums ist $= 4$ Minuten}
\erg{11}{a) Differenz der Raten $d(t) = a'(t) - b'(t) = 36t - 3t^2 - 6t = 30t - 3t^2$; $d'(t) = 30 - 6t = 0$ ⇔ $t = 5$; $d(5) = 75$ cm je Jahr; Ränder: $d(0) = 0$, $d(12) = -72$ – der Unterschied ist nach 5 Jahren mit $75$ cm je Jahr am größten (nicht die Differenz der Höhen untersuchen) \quad b) $d(t) = s'(t) - u'(t) = 0{,}15t^2 + 2 - 3t$; $d'(t) = 0{,}3t - 3 = 0$ ⇔ $t = 10$; $d(10) = -13$, Ränder $d(0) = 2$, $d(15) = -9{,}25$ – nach 10 Sekunden ist der Unterschied mit $13$ m je s am größten (das zweite Fahrzeug ist schneller) \quad c) $d(t) = a'(t) - b'(t) = 3t^2 - 18t + 15$; $d'(t) = 6t - 18 = 0$ ⇔ $t = 3$ mit $d(3) = -12$; Ränder: $d(0) = 15$, $d(6) = 15$ – der Unterschied ist zu Beginn und nach 6 Monaten mit $15$ Tausend je Monat am größten, die innere Extremstelle liefert nur $12$ \quad d) $T(x) - 6 = 50 \cdot \mathrm{e}^{-0{,}02x} = -50 \cdot (-\mathrm{e}^{-0{,}02x}) = -50 \cdot T'(x)$, also $c = -50$ – die Aussage ist wahr; $c$ ist negativ, weil die Temperatur sinkt \quad e) $m(t) = 80 \cdot \mathrm{e}^{-0{,}25t} = -4 \cdot (-20 \cdot \mathrm{e}^{-0{,}25t}) = -4 \cdot m'(t)$, also $c = -4$ – wahr; der Faktor ist negativ, weil die Menge abnimmt \quad f) $L(t) - 100 = -60 \cdot \mathrm{e}^{-0{,}05t} = -20 \cdot (3 \cdot \mathrm{e}^{-0{,}05t}) = -20 \cdot L'(t)$, also $c = -20$ – wahr; der Abstand zur Vollladung ist proportional zur Laderate \quad g) $B(t) - 200 = 500 \cdot \mathrm{e}^{0{,}04t} = 25 \cdot (20 \cdot \mathrm{e}^{0{,}04t}) = 25 \cdot B'(t)$, also $c = 25$ – wahr; hier ist $c$ positiv, weil die Zahl wächst}
